% Independent O/jump construction; not a replacement for a historical source.
\subsection{Effective transfinite computation and delta-one-one definitions}
\label{sec:p04-delta-one-hyperarithmetical}

There are two descriptions of hyperarithmetical sets. One iterates ordinary
oracle jumps along effective ordinal notations. The other requires a set and
its complement to have existential second-order definitions with arithmetic
matrices. The equivalence is a construction theorem: a complementary pair of
definitions can be compiled into a single valid jump oracle and a total
decision program. The following proof supplies that compilation for the
\(\mathcal O\) convention of
Definition~\ref{def:p04-kleene-ordinal-notations}.
It uses neither an assumed Church--Kleene rank bound nor the equality it
proves. Effective grounded recursion and the distinction between ordinal
systems are discussed in \cite[\S\S2A--3E]{Moschovakis2016HYPAuthorManuscript};
the exact \(\mathcal O\) clauses are also displayed in
\cite[\S2.3]{Speranski2017TruthAuthorPreprint}.

\subsubsection{Programs, notations and actual jump sets}

\begin{hblemma}{thm:p04-doh-arithmetic-graphs}{proved=annex:p04.doh.arithmetic-graphs}{Arithmetic formulas for computable numerical graphs}
Every partial computable numerical function has its graph defined by an
arithmetical formula in the standard model with vocabulary
\(0,1,+,\times\), obtained effectively from a machine for the function.
Every computable predicate, finite-word operation and substitution function
therefore has an arithmetic definition. For a total function its graph
has exactly one output at each standard-model input.
\end{hblemma}

\begin{hbconvention}{def:p04-doh-programs}{Arithmetic and oracle coding}
Work over the standard natural numbers in classical ZF set theory. Fix a
computable pairing bijection \(\langle i,x\rangle\) and effective finite
sequence coding. Fix an acceptable numbering \(\varphi_e^B\) of partial
oracle programs: universal simulation and insertion of numerical parameters
are computable operations on program codes, independent of the oracle.
For a set \(B\subseteq\N\), put
\[
 B'=\{e:\varphi_e^B(e)\text{ is defined}\}.
\]
Definedness uses the numerical-output convention: rejection or a physical
halt without a valid numerical output does not define a value. This is
the same Turing-jump convention as in the earlier oracle presentation.
A reduction \(A\leq_T B\) computes the total characteristic function of
\(A\), with values \(0,1\), from oracle \(B\).
For sets \(B_i\), their effective join is
\[
 \bigoplus_i B_i=\{\langle i,x\rangle:x\in B_i\}.
\]
An index compiler constructs a finite program; it need not run the program
or decide whether a prospective ordinal notation is valid. Computable
numerical operations are expressed by their arithmetic graph definitions;
the independent finite-history construction is supplied in
Theorem~\ref{thm:p04-doh-arithmetic-graphs}. It does not depend on the
hyperarithmetical characterization being proved here.
\end{hbconvention}

\begin{hblemma}{thm:p04-doh-program-tools}{proved=annex:p04.doh.program-tools}{Uniform program self-reference and jump simulation}
For every total computable transformation \(f\) of program codes there is
an index \(e\) satisfying \(\varphi_e^B=\varphi_{f(e)}^B\) for every oracle
\(B\). Numerical parameters can remain arguments of this one program.
There are total computable index transformations witnessing
\(B\leq_T B'\), and taking a reduction \(A\leq_T B\) to a reduction
\(A'\leq_T B'\). Their correctness is conditional only on the supplied
reduction actually being total and correct.
\end{hblemma}

\begin{hbdefinition}{def:p04-doh-notations}{The simultaneous notation system}
Write \(\mathsf s(a)=2^a\), and \(\mathsf l(e)=3\cdot5^e\).
Define \(\mathcal O\) and its predecessor relation \(<_{\mathcal O}\)
by the least simultaneous closure under these clauses:
\begin{enumerate}
 \item \(1\in\mathcal O\).
 \item If \(a\in\mathcal O\), then \(\mathsf s(a)\in\mathcal O\),
 and its predecessors are \(a\) and the predecessors of \(a\).
 \item If \(\varphi_e\) is total, all \(a_i=\varphi_e(i)\) belong to
 \(\mathcal O\), and \(a_i<_{\mathcal O}a_{i+1}\) for every \(i\),
 then \(\mathsf l(e)\in\mathcal O\). Its predecessors are exactly the
 \(x\) such that \(x\leq_{\mathcal O}a_i\) for some \(i\), where
 \(\leq_{\mathcal O}\) includes equality.
\end{enumerate}
``Least'' means that no notation or predecessor pair is added except by
iteration of these clauses. A limit clause is an infinite-premise clause;
it is not a computable validity test. Equivalently, take the least fixed
point of the positive operator on tagged notation codes and predecessor
pairs specified by the clauses, retaining already added records at each
step. Set-sized ordinal iteration constructs this fixed point.

The ordinal denotation is
\[
 \nu_{\mathcal O}(1)=0,\qquad
 \nu_{\mathcal O}(\mathsf s(a))=\nu_{\mathcal O}(a)+1,\qquad
 \nu_{\mathcal O}(\mathsf l(e))=\sup_i\nu_{\mathcal O}(\varphi_e(i)).
\]
The associated sets, with the manuscript's empty-base convention, are
\[
 H_1=\varnothing,\qquad H_{\mathsf s(a)}=H_a',\qquad
 H_{\mathsf l(e)}=\bigoplus_i H_{\varphi_e(i)}.
\]
These are definitions on valid notations. Values on invalid codes are not
silently assigned. This is the smaller \(\mathcal O\) system, not the
system which admits every sequence merely increasing in denoted ordinal.
\end{hbdefinition}

\begin{hblemma}{thm:p04-doh-notation-existence}{proved=annex:p04.doh.notation-existence}{Well-defined denotations and hierarchy}
Every valid notation has a unique ordinal denotation and a unique set
\(H_a\). If \(b<_{\mathcal O}a\), then
\(\nu_{\mathcal O}(b)<\nu_{\mathcal O}(a)\).
In particular, repeatedly following a successor predecessor or a selected
limit-sequence term cannot continue infinitely from a valid notation.
\end{hblemma}

\begin{hbdefinition}{def:p04-doh-addition}{A syntactic notation-addition compiler}
The computable compiler \(\operatorname{add}_{\mathcal O}(a,b)\) obeys
\[
\begin{aligned}
 \operatorname{add}_{\mathcal O}(a,1)&=a,\\
 \operatorname{add}_{\mathcal O}(a,\mathsf s(c))
   &=\mathsf s(\operatorname{add}_{\mathcal O}(a,c)),\\
 \operatorname{add}_{\mathcal O}(a,\mathsf l(e))
   &=\mathsf l(\operatorname{index}(i\mapsto
          \operatorname{add}_{\mathcal O}(a,\varphi_e(i)))).
\end{aligned}
\]
The last line outputs an index without running its sequence. On codes of
none of these syntactic forms put \(\operatorname{add}_{\mathcal O}(a,b)=1\).
Program self-reference fixes the sequence index. Successor parsing recurses
only to the smaller natural number \(c\); hence the top-level compiler is
total even on invalid codes. This totality does not assert validity of its
output on those codes.
\end{hbdefinition}

\begin{hblemma}{thm:p04-doh-addition}{proved=annex:p04.doh.addition}{Notation addition and uniform right-summand access}
For \(a,b\in\mathcal O\), the compiled code belongs to \(\mathcal O\) and
\[
 \nu_{\mathcal O}(\operatorname{add}_{\mathcal O}(a,b))
     =\nu_{\mathcal O}(a)+\nu_{\mathcal O}(b).
\]
It preserves the predecessor relation in its second argument. Also
\(a\leq_{\mathcal O}\operatorname{add}_{\mathcal O}(a,b)\), with a strict
inequality when \(b\ne1\).
There is a total computable index function \(r(a,b)\) such that
\[
 \varphi_{r(a,b)}^{H_{\operatorname{add}_{\mathcal O}(a,b)}}
       =\chi_{H_b}.
\]
The reduction does not infer Turing reducibility from ordinal inequality;
its successor simulation and limit-column program are explicitly constructed.
\end{hblemma}

\subsubsection{An effective oracle bound for a well-founded computation}

\begin{hbdefinition}{def:p04-doh-node-bounds}{Node notations for a recursive child relation}
Let \(\mathsf C(z,i)\) be a decidable relation saying that the natural
number \(i\) is a child of node \(z\). Assume that there is no infinite
sequence \(z_0,z_1,\ldots\) with \(\mathsf C(z_j,z_{j+1})\).
Fix a self-referential computable code generator \(h\) as follows. Its
sequence program at node \(z\) computes
\[
\begin{aligned}
 a_{z,0}&=1,\\
 b_{z,i}&=\begin{cases}h(i),&\mathsf C(z,i),\\1,&\text{otherwise},\end{cases}\\
 a_{z,i+1}&=\operatorname{add}_{\mathcal O}
                  (\mathsf s(a_{z,i}),b_{z,i}),\\
 e_z&=\operatorname{index}(i\mapsto a_{z,i}),\\
 h(z)&=\mathsf s(\mathsf s(\mathsf l(e_z))).
\end{aligned}
\]
To output \(h(z)\), the generator only constructs \(e_z\); it does not
evaluate its children. Each individual sequence value uses finitely many
calls to the total top-level generator and notation-addition compiler.
\end{hbdefinition}

\begin{hblemma}{thm:p04-doh-node-bounds}{proved=annex:p04.doh.node-bounds}{Valid node bounds with uniform child access}
The generator is total computable. Every \(h(z)\) is a valid notation.
Put \(J_z=H_{\mathsf l(e_z)}=\bigoplus_i H_{a_{z,i}}\), so
\(H_{h(z)}=J_z''\). For any child \(i\) of \(z\),
\(H_{h(i)}\) is computable from \(J_z\), uniformly in \(z,i\): apply
\(r(\mathsf s(a_{z,i}),h(i))\) using column \(i+1\).
Consequently every child oracle is uniformly computable from \(H_{h(z)}\).
\end{hblemma}

\begin{hbexample}{def:p04-doh-leaf-bound}{A leaf bound is deliberately not minimal}
At a node with no children, every dummy is \(1\), so the cumulative
sequence is \(1,\mathsf s(1),\mathsf s(\mathsf s(1)),\ldots\).
Its limit denotes \(\omega\), and \(h(z)\) denotes \(\omega+2\).
The construction thus gives a uniform valid oracle bound even at a leaf;
it does not claim the least possible ordinal or computational complexity
for a particular recursion. Correctness and effective access, not rank
minimality, are the properties used in the characterization proof.
\end{hbexample}

\begin{hbdefinition}{def:p04-doh-alternation}{An alternating product-tree equation}
Let \(T,U\subseteq\N^{<\N}\) be recursive prefix-closed trees containing
the empty sequence. At a pair \(z=(s,t)\) of equal-length nodes put
\[
 \vartheta(s,t)=1\quad\Longleftrightarrow\quad
 \forall m\bigl[s^\frown m\in T\ \Longrightarrow\
       \exists k\,(t^\frown k\in U\ \land\
                \vartheta(s^\frown m,t^\frown k)=1)\bigr].
\]
The equation assigns a bit by well-founded recursion when its synchronized
product tree is well founded. Its empty-existential case is false, and its
empty-universal case is true. For a recursive family \(T_n,U_n\), use
the numerical code
\[
 \operatorname{node}(n,s,t)=
  1+\langle n,\langle\operatorname{code}(s),\operatorname{code}(t)\rangle\rangle
\]
for a valid equal-length pair, and use \(*=0\) as a common node whose
children are their root pairs. Numbers whose decoded sequences are invalid
or outside the indicated trees have no children. All child and
tree-membership tests are decidable.
\end{hbdefinition}

\begin{hblemma}{thm:p04-doh-alternation}{proved=annex:p04.doh.alternation}{Two-jump evaluation of the alternating equation}
If the common product forest is well founded, the equation has a unique
bit solution. A single effectively compiled oracle program computes
\(\vartheta(s,t)\) from \(H_{h(n,s,t)}\), uniformly in its nodes.
All root bits \(\vartheta(\varnothing,\varnothing)\) are then computed
by one total program from the fixed valid oracle \(H_{h(*)}\).
If \(T\) is well founded and \(U\) has a branch, its root bit is \(1\).
If \(T\) has a branch and \(U\) is well founded, its root bit is \(0\).
\end{hblemma}

\subsubsection{The two inclusions}

\begin{hblemma}{thm:p04-doh-tree-normal-form}{proved=annex:p04.doh.tree-normal-form}{Effective tree normal form}
The function-quantifier convention of
Definition~\ref{def:p04-lightface-analytical-hierarchy} and the
set-quantifier convention used here have effective arithmetical translations
in both directions, preserving \(\Sigma^1_1\) and \(\Pi^1_1\).
Function graphs require arithmetic totality and functionality guards;
set characteristic functions require an arithmetic binary-value guard.
The guard is conjoined under an existential quantifier and made an
antecedent under a universal quantifier.
From a lightface \(\Pi^1_1\) definition of a set \(A\subseteq\N\), one
can effectively construct recursive prefix-closed trees \(T_n\), with
\(n\in A\) exactly when \(T_n\) has no infinite branch. The same
construction converts a lightface existential-set arithmetic definition
into the existence of a branch. No effective selection of a branch is
asserted. This matches the independent tree normal form used in
Theorem~\ref{thm:p11-ktf-tree-normal-form}.
\end{hblemma}

\begin{hbtheorem}{thm:p04-doh-delta-to-hyp}{proved=annex:p04.doh.delta-to-hyp}{Compiling complementary definitions into one jump oracle}
Given a \(\Pi^1_1\) definition of \(A\) and a \(\Pi^1_1\) definition of
\(\N\setminus A\), the construction effectively supplies a natural code
\(a\) and a program index \(d\) such that
\[
 a\in\mathcal O,\qquad \varphi_d^{H_a}=\chi_A.
\]
The promise that the two definitions are complementary is essential:
the compiler need not decide that promise, and correctness or validity is
not asserted for an arbitrary incompatible pair of definitions.
Thus every lightface \(\Delta^1_1\) subset of \(\N\) is hyperarithmetical.
\end{hbtheorem}

\begin{hbtheorem}{thm:p04-doh-hyp-to-delta}{proved=annex:p04.doh.hyp-to-delta}{Arithmetic specification of a unique jump hierarchy}
For each fixed \(a\in\mathcal O\), both \(H_a\) and its complement have
lightface \(\Sigma^1_1\) definitions. This holds uniformly by a finite
formula with the notation code as a numerical parameter, with no claim of
validity on invalid codes. The same is true of every set Turing reducible
to \(H_a\). Hence \(\mathsf{HYP}\subseteq\Delta^1_1\).
\end{hbtheorem}

\begin{hbcorollary}{thm:p04-doh-characterization}{proved=annex:p04.doh.characterization}{The exact hyperarithmetical characterization}
With the actual notation and jump conventions above,
\[
 \{A\subseteq\N:\exists a\in\mathcal O\ (A\leq_T H_a)\}
       =\Delta^1_1.
\]
Numerical parameters are harmless: each fixed finite tuple can be inserted
into a lightface definition or a program. Arbitrary set parameters are not
implicitly allowed. Relative analytical pointclasses and relative ordinal
notation systems require their own explicitly fixed oracle conventions;
they are not the unrelativized \(\mathcal O\) with an arbitrary set
smuggled into its limit sequences.
\end{hbcorollary}

The compiler proof is independent of the finite-level arithmetical hierarchy
and of arithmetic-truth closure ordinals. In particular, a total monadic
extension already shown to be lightface \(\Delta^1_1\) can now be identified
as hyperarithmetical without making that classification an assumption of
the truth construction.
